% !TeX root = ../教材题库_全部条目_转写工作稿.tex
% Source: XJ-B1-C02-EX-P039-Y0303
% 第2章 / 2.1 相等关系与不等关系 / 教材 p.39 / PDF 第46页
% OCR未校对：公式与文字均可能有误；未完成数学TeX转写；禁止直接发给学生
% 原题及相邻原解截图（核对用）：../题目原图/XJ-B1-C02-EX-P039-Y0303-01.webp
% 本文件只在转写完成后作为题干使用；OCR 条目尚不是可靠题干。
\begin{bankitem}{2.1 例题 例题5}{第2章 / 2.1 相等关系与不等关系 / 教材 p.39 / PDF 第46页}{XJ-B1-C02-EX-P039-Y0303}{OCR未校对：公式与文字均可能有误；未完成数学TeX转写；禁止直接发给学生}
\begin{ocrdraft}
[教材p.39 / PDF第46页]
\par
00设。乃为正数，证明下列不等式：
\par
（1）磊十一冫2；
\par
（2）一十莖冫2．
\par
证明（D因为a，一均为正数，由基本不等式，得
\par
1
\par
a十一冫2磊
\par
1
\par
当且仅当a==
\par
，即=1时等号成立，所以原不等式成立．
\par
（2）因为乃为正数，所以
\par
莖也为正数，由基本不等式，得
\par
一十莖冫2
\par
当且仅当
\par
，即一b时等号成立，所以原不等式成立．
\end{ocrdraft}
\end{bankitem}
